\documentclass{amsart}
\usepackage[UTF8,fontset=windows]{ctex}
\usepackage[a4paper,margin=24mm]{geometry}
% For dealing with references we use the comment environment
\usepackage{verbatim}
\newenvironment{reference}{\comment}{\endcomment}
%\newenvironment{reference}{}{}
\newenvironment{slogan}{\comment}{\endcomment}
\newenvironment{history}{\comment}{\endcomment}

% For commutative diagrams we use Xy-pic
\usepackage[all]{xy}

% We use 2cell for 2-commutative diagrams.
\xyoption{2cell}
\UseAllTwocells

% We use multicol for the list of chapters between chapters
\usepackage{multicol}

% This is generally recommended for better output
\usepackage{lmodern}
\usepackage[T1]{fontenc}

% For cross-file-references

% Package for hypertext links:
\usepackage{hyperref}

% For any local file, say "hello.tex" you want to link to please

% Theorem environments.
%
\theoremstyle{plain}
\newtheorem{theorem}[subsection]{Theorem}
\newtheorem{proposition}[subsection]{Proposition}
\newtheorem{lemma}[subsection]{Lemma}

\theoremstyle{definition}
\newtheorem{definition}[subsection]{Definition}
\newtheorem{example}[subsection]{Example}
\newtheorem{exercise}[subsection]{Exercise}
\newtheorem{situation}[subsection]{Situation}

\theoremstyle{remark}
\newtheorem{remark}[subsection]{Remark}
\newtheorem{remarks}[subsection]{Remarks}

\numberwithin{equation}{subsection}

% Macros
%
\def\lim{\mathop{\mathrm{lim}}\nolimits}
\def\colim{\mathop{\mathrm{colim}}\nolimits}
\def\Spec{\mathop{\mathrm{Spec}}}
\def\Hom{\mathop{\mathrm{Hom}}\nolimits}
\def\Ext{\mathop{\mathrm{Ext}}\nolimits}
\def\SheafHom{\mathop{\mathcal{H}\!\mathit{om}}\nolimits}
\def\SheafExt{\mathop{\mathcal{E}\!\mathit{xt}}\nolimits}
\def\Sch{\mathit{Sch}}
\def\Mor{\mathop{\mathrm{Mor}}\nolimits}
\def\Ob{\mathop{\mathrm{Ob}}\nolimits}
\def\Sh{\mathop{\mathit{Sh}}\nolimits}
\def\NL{\mathop{N\!L}\nolimits}
\def\CH{\mathop{\mathrm{CH}}\nolimits}
\def\proetale{{pro\text{-}\acute{e}tale}}
\def\etale{{\acute{e}tale}}
\def\QCoh{\mathit{QCoh}}
\def\Ker{\mathop{\mathrm{Ker}}}
\def\Im{\mathop{\mathrm{Im}}}
\def\Coker{\mathop{\mathrm{Coker}}}
\def\Coim{\mathop{\mathrm{Coim}}}

% Boxtimes
%
\DeclareMathSymbol{\boxtimes}{\mathbin}{AMSa}{"02}

%
% Macros for moduli stacks/spaces
%
\def\QCohstack{\mathcal{QC}\!\mathit{oh}}
\def\Cohstack{\mathcal{C}\!\mathit{oh}}
\def\Spacesstack{\mathcal{S}\!\mathit{paces}}
\def\Quotfunctor{\mathrm{Quot}}
\def\Hilbfunctor{\mathrm{Hilb}}
\def\Curvesstack{\mathcal{C}\!\mathit{urves}}
\def\Polarizedstack{\mathcal{P}\!\mathit{olarized}}
\def\Complexesstack{\mathcal{C}\!\mathit{omplexes}}
% \Pic is the operator that assigns to X its picard group, usage \Pic(X)
% \Picardstack_{X/B} denotes the Picard stack of X over B
% \Picardfunctor_{X/B} denotes the Picard functor of X over B
\def\Pic{\mathop{\mathrm{Pic}}\nolimits}
\def\Picardstack{\mathcal{P}\!\mathit{ic}}
\def\Picardfunctor{\mathrm{Pic}}
\def\Deformationcategory{\mathcal{D}\!\mathit{ef}}

\hypersetup{hidelinks}
\widowpenalty=10000\clubpenalty=10000\displaywidowpenalty=10000
\begin{document}
\title[Stacks Project, Tag 0FLA]{局部拟有限 syntomic 态射的微分形式比较映射}
\maketitle
\noindent\textbf{Counting scope.} The primary problem is Tag 0FLA, marked below. All preceding numbered statements are supporting author context and are not counted as additional problems. This is a local pilot with a temporary ID, not production admission.


\noindent Copyright (C) 2005--2025 Johan de Jong. Stacks Project authors. Source: \href{https://stacks.math.columbia.edu/tag/0FLA}{Tag 0FLA}.

\noindent Difficulty H2: 本题构造对所有局部拟有限syntomic态射自然的微分形式到det(NL)扭曲微分形式的映射。允许局部模型及基变换引理后，仍须用纤维积比较两个nice模型，证明所得局部映射与模型选择无关并能粘合。核心局部构造有明确前置输入，而多模型一致性与全局化是本题承担的实质任务，判H2。.

\noindent History: excerpt from revision local cached snapshot; exact chapter hashes; revision unknown, derham.tex, lines 3520--3764. Only reference presentation was adapted. Licensed under GNU FDL 1.2 or later; no invariant sections or cover texts. See accompanying COPYING.

\section{Comparing sheaves of differential forms}
\label{section-quasi-finite-syntomic}

\noindent
The goal of this section is to construct for every locally quasi-finite
syntomic morphism of schemes $f : Y \to X$ maps
$$
c^p_{Y/X} :
\Omega^p_{Y/\mathbf{Z}}
\longrightarrow
f^*\Omega^p_{X/\mathbf{Z}} \otimes_{\mathcal{O}_Y} \det(\NL_{Y/X})
$$
for all $p \geq 0$ satisfying the following properties
\begin{enumerate}
\item
\label{item-degree-zero}
$c_{Y/X}^0(1) = \delta(\NL_{Y/X})$ see
Discriminants, Section \href{https://stacks.math.columbia.edu/tag/0FKB}{section tate map (Tag 0FKB)},
\item
\label{item-multiplicative}
$c_{Y/X}^{q + p}(\omega \wedge \eta) = \omega \wedge c_{Y/X}^p(\eta)$
for local sections $\omega$ of $f^*\Omega^q_{X/\mathbf{Z}}$
and $\eta$ of $\Omega^p_{Y/\mathbf{Z}}$,
\item
\label{item-base-change}
the formation of $c^p_{Y/X}$ is compatible with restriction to opens
and with base change
(as formulated in Lemma \ref{lemma-base-change-Garel-upstairs}).
\end{enumerate}
Note that these conditions imply that the maps $c^p_{Y/X}$ are
$\mathcal{O}_Y$-linear and that the composition with
$f^*\Omega^p_{X/\mathbf{Z}} \to \Omega^p_{Y/\mathbf{Z}}$
is multiplication by $\delta(\NL_{Y/X})$.
We will in fact show that there exists a unique collection of operators
$c^p_{Y/X}$ with properties (1), (2), (3). This result will be applied in
Section \href{https://stacks.math.columbia.edu/tag/0FK6}{section trace (Tag 0FK6)} to the construction of the trace map on
the de Rham complexes if $f$ is finite.



\begin{remark}
\label{remark-local-description}
Let $A$ be a ring. Let $P = A[x_1, \ldots, x_n]$. Let
$f_1, \ldots, f_n \in P$ and set $B = P/(f_1, \ldots, f_n)$.
Assume $A \to B$ is quasi-finite. Then
$B$ is a relative global complete intersection over $A$ (Algebra, Definition
\ref{algebra-definition-relative-global-complete-intersection}) and
$(f_1, \ldots, f_n)/(f_1, \ldots, f_n)^2$ is free with generators
the classes $\overline{f}_i$ by Algebra, Lemma
\href{https://stacks.math.columbia.edu/tag/00SV}{lemma relative global complete intersection conormal (Tag 00SV)}.
Consider the following diagram
$$
\xymatrix{
\Omega_{A/\mathbf{Z}} \otimes_A B \ar[r] &
\Omega_{P/\mathbf{Z}} \otimes_P B \ar[r] &
\Omega_{P/A} \otimes_P B \\
&
(f_1, \ldots, f_n)/(f_1, \ldots, f_n)^2 \ar[u] \ar@{=}[r] &
(f_1, \ldots, f_n)/(f_1, \ldots, f_n)^2 \ar[u]
}
$$
The right column represents $\NL_{B/A}$ in $D(B)$ hence has cohomology
$\Omega_{B/A}$ in degree $0$. The top row is the split short exact sequence
$0 \to \Omega_{A/\mathbf{Z}} \otimes_A B \to
\Omega_{P/\mathbf{Z}} \otimes_P B \to \Omega_{P/A} \otimes_P B \to 0$.
The middle column has cohomology $\Omega_{B/\mathbf{Z}}$ in degree $0$
by Algebra, Lemma \href{https://stacks.math.columbia.edu/tag/00RU}{lemma differential seq (Tag 00RU)}.
Thus by Lemma \href{https://stacks.math.columbia.edu/tag/0FL8}{lemma funny map (Tag 0FL8)} we obtain canonical $B$-module maps
$$
\Omega^p_{B/\mathbf{Z}} \longrightarrow
\Omega^p_{A/\mathbf{Z}} \otimes_A \det(\NL_{B/A})
$$
for $p \geq 0$ satisfying properties (1) and (2) and in particular
satisfy the property that the composition with
$\Omega^p_{A/\mathbf{Z}} \to \Omega^p_{B/\mathbf{Z}}$
is multiplication by $\delta(\NL_{B/A})$.
\end{remark}

\begin{lemma}
\label{lemma-base-change-Garel-upstairs}
Consider a commutative diagram
$$
\xymatrix{
Y' \ar[d]_{f'} \ar[r]_b & Y \ar[d]^f \\
X' \ar[r]^a & X
}
$$
of schemes which induces an isomorphism of $Y'$ with an open subscheme of
$X' \times_X Y$. Assume $f$ is locally quasi-finite and syntomic and
assume given maps $c^p_{Y/X} : \Omega^p_{Y/\mathbf{Z}} \to
f^*\Omega^p_{X/\mathbf{Z}} \otimes_{\mathcal{O}_Y} \det(\NL_{Y/X})$
for $p \geq 0$ satisfying
(\ref{item-degree-zero}) and (\ref{item-multiplicative}).
Then there is at most one collection of maps
$c^p_{Y'/X'} : \Omega^p_{Y'/\mathbf{Z}} \to
(f')^*\Omega^p_{X'/\mathbf{Z}} \otimes_{\mathcal{O}_{Y'}} \det(\NL_{Y'/X'})$
for $p \geq 0$
satisfying (\ref{item-degree-zero}) and (\ref{item-multiplicative})
such that the diagrams
$$
\xymatrix{
b^*\Omega^p_{Y/\mathbf{Z}} \ar[rr]_-{b^*c^p_{Y/X}} \ar[d] & &
b^*(f^*\Omega^p_{X/\mathbf{Z}} \otimes_{\mathcal{O}_Y} \det(\NL_{Y/X}))
\ar@{=}[r] &
(f')^*a^*\Omega^p_{X/\mathbf{Z}} \otimes_{\mathcal{O}_Y} b^*\det(\NL_{Y/X}))
\ar[d] \\
\Omega^p_{Y'/\mathbf{Z}} \ar[rrr]^-{c^p_{Y'/X'}} & & &
(f')^*\Omega^p_{X'/\mathbf{Z}} \otimes_{\mathcal{O}_{Y'}} \det(\NL_{Y'/X'})
}
$$
commute for all $p \geq 0$. Here the vertical arrows use the maps
$b^*\Omega^p_{Y/\mathbf{Z}} \to \Omega^p_{Y'/\mathbf{Z}}$ and
$a^*\Omega^p_{X/\mathbf{Z}} \to \Omega^p_{X'/\mathbf{Z}}$
of Section \href{https://stacks.math.columbia.edu/tag/07HX}{section de rham complex (Tag 07HX)}
as well as the identification $b^*\det(\NL_{Y/X}) = \det(\NL_{Y'/X'})$ of
Discriminants, Section \href{https://stacks.math.columbia.edu/tag/0FKB}{section tate map (Tag 0FKB)}.
\end{lemma}

\begin{proof}
The map
$$
(f')^*\Omega_{X'/\mathbf{Z}} \oplus b^*\Omega_{Y/\mathbf{Z}}
\longrightarrow
\Omega_{Y'/\mathbf{Z}}
$$
is surjective because $Y'$ is an open subscheme of the fibre product;
the corresponding algebra statement is that $\Omega_{B \otimes_A C/\mathbf{Z}}$
is generated as a module by the images of $\Omega_{B/\mathbf{Z}}$ and
$\Omega_{C/\mathbf{Z}}$. Thus for all $p$ the map
$$
\bigoplus\nolimits_{i = 0, \ldots, p}
(f')^*\Omega^i_{X'/\mathbf{Z}} \otimes
b^*\Omega^{p - i}_{Y/\mathbf{Z}}
\longrightarrow
\Omega^p_{Y'/\mathbf{Z}}
$$
is surjective. Conditions (\ref{item-degree-zero}) and
(\ref{item-multiplicative}) combined with the commutativity
of the diagram in the lemma, implies that the map $c^p_{Y'/X'}$
is determined (but existence is not immediate).
\end{proof}


\begin{lemma}
\label{lemma-existence-base-change-Garel-upstairs}
Consider a commutative diagram
$$
\xymatrix{
Y' \ar[d]_{f'} \ar[r]_b & Y \ar[d]^f \\
X' \ar[r]^a & X
}
$$
of schemes which induces an isomorphism of $Y'$ with an open subscheme of
$X' \times_X Y$. Assume $f$ is locally quasi-finite and syntomic. Then for
every $y' \in Y'$ we can find opens $V' \subset Y'$,
$V \subset Y$, $U' \subset X$, $U \subset X$ with $y' \in V'$, with
$f'(V') \subset U'$, $b(V') \subset V$, $a(U') \subset U$,
$f(V) \subset U$, and such that for $p \geq 0$ there are maps
$c^p_{V/U}$ and $c^p_{V'/U'}$ satisfying (\ref{item-degree-zero})
and (\ref{item-multiplicative}) which are compatible with the diagram
$$
\xymatrix{
V' \ar[d] \ar[r]_b & V \ar[d] \\
U' \ar[r] & U
}
$$
in the sense explained in Lemma \ref{lemma-base-change-Garel-upstairs}.
\end{lemma}

\begin{proof}
It is clear that we may replace $Y'$ by $X' \times_X Y$ in order to prove this.
Pick affine opens $V \subset Y$ and $U \subset X$ with $V$ containing
the image of $y'$ as in
Discriminants, Lemma \href{https://stacks.math.columbia.edu/tag/0BWE}{lemma syntomic quasi finite (Tag 0BWE)} part (5).
Pick an affine open $U' \subset X'$ containing the image of $y'$
and mapping into $U$. After replacing $X, Y, X', Y'$ by
$U, V, U', U' \times_U V$ we may assume we are in the situation
described in the next paragraph.

\medskip\noindent
Assume that $X' = \Spec(A')$, $X = \Spec(A)$, $Y = \Spec(B)$, and
$Y' = \Spec(A' \otimes_A B)$ with $B = A[x_1, \ldots, x_n]/(f_1, \ldots, f_n)$
a relative global complete intersection over $A$. Then also
$B' = A'[x_1, \ldots, x_n]/(f'_1, \ldots, f'_n)$ is a relative
global complete intersection over $A'$ where $f'_i$ is the image of $f_i$
in $A'[x_1, \ldots, x_n]$, see Algebra, Lemma
\href{https://stacks.math.columbia.edu/tag/00SS}{lemma base change relative global complete intersection (Tag 00SS)}.
The construction in Remark \ref{remark-local-description}
provides us with the maps $c^p_{V/U}$ and $c^p_{V'/U'}$.
These maps are compatible by Lemma \href{https://stacks.math.columbia.edu/tag/0H9F}{lemma funny map functorial (Tag 0H9F)}
and the compatibility of the presentations of $B$ and $B'$ over
$A$ and $A'$. Some details omitted.
\end{proof}


\begin{example}
\label{discriminant-example-universal-quasi-finite-syntomic}
Let $n \geq 1$ and $d \geq 1$ be integers. Let $T$ be the set of
multi-indices $E = (e_1, \ldots, e_n)$ with $e_i \geq 0$ and
$\sum e_i \leq d$. Consider the ring
$$
A = \mathbf{Z}[a_{i, E} ; 1 \leq i \leq n, E \in T]
$$
In $A[x_1, \ldots, x_n]$ consider the elements
$f_i = \sum_{E \in T} a_{i, E} x^E$ where $x^E = x_1^{e_1} \ldots x_n^{e_n}$
as is customary. Consider the $A$-algebra
$$
B = A[x_1, \ldots, x_n]/(f_1, \ldots, f_n)
$$
Denote $X_{n, d} = \Spec(A)$ and let $Y_{n, d} \subset \Spec(B)$
be the maximal open subscheme such that the restriction of the
morphism $\Spec(B) \to \Spec(A) = X_{n, d}$ is quasi-finite, see
Algebra, Lemma \href{https://stacks.math.columbia.edu/tag/00QA}{lemma quasi finite open (Tag 00QA)}.
\end{example}

\begin{definition}
\label{algebra-definition-relative-global-complete-intersection}
Let $R \to S$ be a ring map. We say that $R \to S$ is
a {\it relative global complete intersection} if there exists
a presentation $S = R[x_1, \ldots, x_n]/(f_1, \ldots, f_c)$ and
every nonempty fibre of $\Spec(S) \to \Spec(R)$ has dimension $n - c$.
We will say ``let $S = R[x_1, \ldots, x_n]/(f_1, \ldots, f_c)$ be a relative
global complete intersection'' to indicate this situation.
\end{definition}

\section*{Primary problem: Tag 0FLA}
\begin{lemma}
\label{lemma-Garel-upstairs}
There exists a unique rule that to every locally quasi-finite syntomic
morphism of schemes $f : Y \to X$ assigns $\mathcal{O}_Y$-module maps
$$
c^p_{Y/X} :
\Omega^p_{Y/\mathbf{Z}}
\longrightarrow
f^*\Omega^p_{X/\mathbf{Z}} \otimes_{\mathcal{O}_Y} \det(\NL_{Y/X})
$$
for $p \geq 0$ satisfying (\ref{item-degree-zero}), (\ref{item-multiplicative}),
and (\ref{item-base-change}). In particular, the composition of
$c^p_{Y/X}$ with $f^*\Omega^p_{X/\mathbf{Z}} \to \Omega^p_{Y/\mathbf{Z}}$
is multiplication by $\delta(\NL_{Y/X})$.
\end{lemma}

\begin{proof}
This proof is very similar to the proof of
Discriminants, Proposition \href{https://stacks.math.columbia.edu/tag/0FKD}{proposition tate map (Tag 0FKD)}
and we suggest the reader look at that proof first.

\medskip\noindent
Let us reformulate the statement. Consider the category
$\mathcal{C}$ whose objects, denoted $Y/X$, are locally quasi-finite syntomic
morphism $f : Y \to X$ of schemes and whose morphisms
$b/a : Y'/X' \to Y/X$ are commutative diagrams
$$
\xymatrix{
Y' \ar[d]_{f'} \ar[r]_b & Y \ar[d]^f \\
X' \ar[r]^a & X
}
$$
which induce an isomorphism of $Y'$ with an open subscheme of
$X' \times_X Y$. The lemma means that for every object
$Y/X$ of $\mathcal{C}$ we have maps $c^p_{Y/X}$, $p \geq 0$
with properties (\ref{item-degree-zero}) and (\ref{item-multiplicative}),
and for every morphism $b/a : Y'/X' \to Y/X$ of $\mathcal{C}$ we have
that $c^p_{Y/X}$ and $c^p_{Y'/X'}$ are compatible as in
Lemma \ref{lemma-base-change-Garel-upstairs}.

\medskip\noindent
Given $Y/X$ in $\mathcal{C}$ and $y \in Y$ we can find
an affine open $V \subset Y$ and $U \subset X$ with $f(V) \subset U$
such that there exists some maps
$$
\Omega^p_{Y/\mathbf{Z}}|_V
\longrightarrow
\left(
f^*\Omega^p_{X/\mathbf{Z}} \otimes_{\mathcal{O}_Y} \det(\NL_{Y/X})
\right)|_V
$$
with properties (\ref{item-degree-zero}) and (\ref{item-multiplicative}).
This follows from picking affine opens as in
Discriminants, Lemma \href{https://stacks.math.columbia.edu/tag/0BWE}{lemma syntomic quasi finite (Tag 0BWE)} part (5)
and then using the construction in Remark \ref{remark-local-description}.

\medskip\noindent
Note that the \'etale locus of $f$ is exactly the set of points where
$\delta(\NL_{Y/X})$ does not vanish, see discussion in
Discriminants, Section \href{https://stacks.math.columbia.edu/tag/0FKB}{section tate map (Tag 0FKB)}.
In particular
$f^*\Omega^p_{X/\mathbf{Z}} \to \Omega^p_{Y/\mathbf{Z}}$
becomes an isomorphism after inverting $\delta(\NL_{Y/X})$.
We conclude that if $\Omega^p_{X/\mathbf{Z}}$ is finite locally free and
the annihilator of $\delta(\NL_{Y/X})$ in $\mathcal{O}_Y$ is zero, then
the local maps constructed in the previous paragraph
are unique and automatically glue!

\medskip\noindent
Let $\mathcal{C}_{nice} \subset \mathcal{C}$ denote the full subcategory
of $Y/X$ such that
\begin{enumerate}
\item $\Omega_{X/\mathbf{Z}}$ is locally free, and
\item the annihilator of $\delta(\NL_{Y/X})$ in $\mathcal{O}_Y$ is zero.
\end{enumerate}
By the remarks in the previous paragraph, we see that for any
object $Y/X$ of $\mathcal{C}_{nice}$ we have unique maps
$c^p_{Y/X}$, $p \geq 0$ satisfying conditions (\ref{item-degree-zero})
and (\ref{item-multiplicative}). If $b/a : Y'/X' \to Y/X$
is a morphism of $\mathcal{C}_{nice}$, then the maps
$c^p_{Y/X}$ and $c^p_{Y'/X'}$ are compatible as in
Lemma \ref{lemma-base-change-Garel-upstairs}: namely,
locally there do exist compatible maps by
Lemma \ref{lemma-existence-base-change-Garel-upstairs}
and the uniqueness just mentioned shows these maps agree
with $c^p_{Y/X}$ and $c^p_{Y'/X'}$.
In other words, we have solved the problem
on the full subcategory $\mathcal{C}_{nice}$. For $Y/X$ in $\mathcal{C}_{nice}$
we continue to denote $c^p_{Y/X}$ the solution we've just found.

\medskip\noindent
In fact, more generally, suppose we have a morphism $b/a : Y'/X' \to Y/X$
of $\mathcal{C}$ such that $Y/X$ is an object of $\mathcal{C}_{nice}$.
The same argument, but this time using the uniqueness in
Lemma \ref{lemma-base-change-Garel-upstairs}, tells us there are
maps $c^p_{Y'/X'}$, $p \geq 0$  satisfying conditions (\ref{item-degree-zero})
and (\ref{item-multiplicative}) compatible with the already constructed
maps $c^p_{Y/X}$. Let us call this the base change
$(b/a)^*c^p_{Y/X}$.

\medskip\noindent
Consider morphisms
$$
Y_1/X_1 \xleftarrow{b_1/a_1} Y/X \xrightarrow{b_2/a_2} Y_2/X_2
$$
in $\mathcal{C}$ such that $Y_1/X_1$ and $Y_2/X_2$ are objects
of $\mathcal{C}_{nice}$.
{\bf Claim.} The two base changes $(b_i/a_i)^*c^p_{Y_i/X_i}$, $i = 1, 2$
are equal. We will first show that the claim implies the lemma
and then we will prove the claim.

\medskip\noindent
Let $d, n \geq 1$ and consider the locally
quasi-finite syntomic morphism $Y_{n, d} \to X_{n, d}$
constructed in Discriminants, Example
\ref{discriminant-example-universal-quasi-finite-syntomic}.
Then $Y_{n, d}$ and $Y_{n, d}$ are irreducible schemes of finite type and
smooth over $\mathbf{Z}$. Namely, $X_{n, d}$ is a spectrum of a
polynomial ring over $\mathbf{Z}$ and $Y_{n, d}$ is an open subscheme
of such. The morphism $Y_{n, d} \to X_{n, d}$ is locally quasi-finite syntomic
and \'etale over a dense open, see Discriminants, Lemma
\href{https://stacks.math.columbia.edu/tag/0FK9}{lemma universal quasi finite syntomic etale (Tag 0FK9)}.
Thus $\delta(\NL_{Y_{n, d}/X_{n, d}})$ is nonzero: for example we have
the local description of $\delta(\NL_{Y/X})$ in
Discriminants, Remark \href{https://stacks.math.columbia.edu/tag/0FKC}{remark local description delta (Tag 0FKC)}
and we have the local description of \'etale morphisms in
Morphisms, Lemma \href{https://stacks.math.columbia.edu/tag/02GU}{lemma etale at point (Tag 02GU)} part (8).
Now a nonzero section of an invertible module over an irreducible
regular scheme has vanishing annihilator. Thus
$Y_{n, d}/X_{n, d}$ is an object of $\mathcal{C}_{nice}$.

\medskip\noindent
Let $Y/X$ be an arbitrary object of $\mathcal{C}$. Let $y \in Y$.
By Discriminants, Lemma \href{https://stacks.math.columbia.edu/tag/0FKA}{lemma locally comes from universal (Tag 0FKA)}
we can find $n, d \geq 1$ and morphisms
$$
Y/X \leftarrow V/U \xrightarrow{b/a} Y_{n, d}/X_{n, d}
$$
of $\mathcal{C}$ such that $V \subset Y$ and $U \subset X$ are open.
Then we have the base change $c^p_{V/U} = (b/a)^*c^p_{Y_{n, d}/X_{n, d}}$.
The claim guarantees these locally constructed maps $c^p_{V/U}$ glue!
Thus we get a well defined global maps $c^p_{Y/X}$ with properties
(\ref{item-degree-zero}) and (\ref{item-multiplicative}).
Finally, let $b/a : Y'/X' \to Y/X$ be an arbitray morphism of $\mathcal{C}$.
We have the maps $c^p_{Y'/X'}$, $p \geq 0$ and $c^p_{Y/X}$, $p \geq 0$
constructed in this paragraph. To check they are compatible as in
Lemma \ref{lemma-base-change-Garel-upstairs} we may work locally
on $Y'/X'$ and $Y/X$. Thus we may assume that there exists a morphism
$Y/X \to Y_{n, d}/X_{n, d}$. By construction, both the family
$c^p_{Y'/X'}$ and the family $c^p_{X/Y}$ are compatible with
the map towards $Y_{n, d}/X_{n, d}$. It follows easily from this
(and the uniqueness in Lemma \ref{lemma-base-change-Garel-upstairs})
that $c^p_{Y'/X'}$ is compatible with $c^p_{Y/X}$.
Thus it remains to prove the claim.

\medskip\noindent
In the rest of the proof we prove the claim. We may pick a point
$y \in Y$ and prove the maps agree in an open neighbourhood of $y$.
Thus we may replace $Y_1$, $Y_2$ by open neighbourhoods of the
image of $y$ in $Y_1$ and $Y_2$. Therefore we may assume
$Y, X, Y_1, X_1, Y_2, X_2$ are affine, say they are the spectra
of rings $B, A, B_1, A_1, B_2, A_2$. Picture
$$
\xymatrix{
B_1 \ar[r] & B & B_2 \ar[l] \\
A_1 \ar[u] \ar[r] & A \ar[u] & A_2 \ar[l] \ar[u]
}
$$
By assumption the spectrum of $B$ is an affine open of
both the spectrum of $A \otimes_{A_1} B_1$ and $A \otimes_{A_2} B_2$.
Shrinking more we may assume there exist elements
$g_i \in A \otimes_{A_i} B_i$ such that our maps give isomorphisms
$(A \otimes_{A_i} B_i)_{g_i} = B$, see
Properties, Lemma \href{https://stacks.math.columbia.edu/tag/0H9B}{lemma standard open two affines (Tag 0H9B)}.
Let $x_\alpha$, $y_\beta$ be a sufficiently large collection of
variables such that we may choose surjections
$$
A'_1 = A_1[x_\alpha] \to A
\quad\text{and}\quad
A'_2 = A_2[y_\beta] \to A
$$
of $A_1$ and $A_2$-algebras. Then we can choose lifts
$h_i \in A'_i \otimes_{A_i} B_i$ of $g_i \in A \otimes_{A_i} B_i$
and we consider the diagram
$$
\xymatrix{
(A'_1 \otimes_{A_1} B_1)_{h_1} \ar[r] & B &
(A'_2 \otimes_{A_1} B_2)_{h_2} \ar[l] \\
A'_1 \ar[u] \ar[r] & A \ar[u] & A'_2 \ar[l] \ar[u]
}
$$
By construction the two squares are cocartesian. Next, we consider the
ring map
$$
A' = A'_1 \times_A A'_2 \longrightarrow
B' =
(A'_1 \otimes_{A_1} B_1)_{h_1} \times_B
(A'_2 \otimes_{A_1} B_2)_{h_2}
$$
By
More on Algebra, Lemma \href{https://stacks.math.columbia.edu/tag/07RU}{lemma module over fibre product (Tag 07RU)}
we have $A'_1 \otimes_{A'} B' = (A'_1 \otimes_{A_1} B_1)_{h_1}$
and $A'_2 \otimes_{A'} B' = (A'_2 \otimes_{A_2} B_2)_{h_2}$. In particular
the fibres of the morphism $Y' = \Spec(B') \to \Spec(A') = X'$ are
open subschemes of base changes of the fibres of the maps $Y_i \to X_i$, hence
finite and local complete intersections (see
Discriminants, Lemma \href{https://stacks.math.columbia.edu/tag/0BWE}{lemma syntomic quasi finite (Tag 0BWE)}).
By More on Algebra, Lemma
\href{https://stacks.math.columbia.edu/tag/08KQ}{lemma properties algebras over fibre product (Tag 08KQ)}
the ring map $A' \to B'$ is flat and of finite presentation.
Thus by Discriminants, Lemma
\href{https://stacks.math.columbia.edu/tag/0BWE}{lemma syntomic quasi finite (Tag 0BWE)} part (3)
we conclude that $Y'/X'$ is an object of $\mathcal{C}$.
Consider now the commutative diagram
$$
\xymatrix{
& & Y/X \ar[ld] \ar@/_2pc/[lldd]_{b_1/a_1} \ar[rd] \ar@/^2pc/[rrdd]^{b_2/a_2} \\
& Y'_1/X'_1 \ar[rd] \ar[ld] & & Y'_2/X'_2 \ar[ld] \ar[rd] \\
Y_1/X_1 & & Y'/X' & & Y_2/X_2
}
$$
where $Y'_i/X'_i$ corresponds to $A'_i \to (A'_i \otimes_{A_i} B_i)_{h_i}$.
Note that $Y'_i/X'_i$ is an object of $\mathcal{C}_{nice}$ as it
obtained by taking an open subscheme of the product of an
infinite dimensional affine space with $Y_i/X_i$; small detail omitted.
In particular, the pullback of $c^p_{Y_i/X_i}$ via $(b_i/a_i)$ is the same
as the pullback of $c^p_{Y'_i/X'_i}$ via $Y/X \to Y'_i/X'_i$.
Now we would be done if $Y'/X'$ is an object of $\mathcal{C}_{nice}$,
but this is almost never the case. Namely, then pulling back $c^p_{Y'/X'}$
around the two sides of the square, we would obtain the desired conclusion.
To get around the problem that $Y'/X'$ is not in $\mathcal{C}_{nice}$
we note the arguments above show that, after possibly shrinking all
of the schemes $X, Y, X'_1, Y'_1, X'_2, Y'_2, X', Y'$ we can find some
$n, d \geq 1$, and extend the diagram like so:
$$
\xymatrix{
& Y/X \ar[ld] \ar[rd] \\
Y'_1/X'_1 \ar[rd] & & Y'_2/X'_2 \ar[ld] \\
& Y'/X' \ar[d] \\
& Y_{n, d}/X_{n, d}
}
$$
and then we can use the already given argument by pulling
back from $c^p_{Y_{n, d}/X_{n, d}}$. This finishes the proof.
\end{proof}
\end{document}
